\documentclass[10pt]{extarticle}
\usepackage[a4paper,margin=1.5cm]{geometry}
\usepackage[T1]{fontenc}
\usepackage{amsmath,amssymb}
\usepackage{graphicx}
\usepackage[export]{adjustbox}
\usepackage{multicol}
\usepackage{enumitem}
\usepackage{booktabs,array,tabularx}
\setlength{\columnsep}{0.7cm}
\setlength{\parindent}{0pt}
\setlength{\parskip}{3pt}
\pagestyle{empty}
\makeatletter
% figure registry (figdefs.tex is generated by the extract step)
\newcommand{\qfigdef}[4]{\@namedef{qf@file@#1}{#2}\@namedef{qf@w@#1}{#3}\@namedef{qf@wide@#1}{#4}}
\newcommand{\qf@check}[1]{\@ifundefined{qf@file@#1}{\PackageError{score2}{Unknown figure id #1}{}}{}}
\newcommand{\qf@img}[4][]{\includegraphics[width=#2\dimexpr\@nameuse{qf@w@#4}pt\relax,max width=#3,max height=0.45\textheight,#1]{\@nameuse{qf@file@#4}}}
% Full-width questions: finalize wraps every question that contains a wide figure in \qfullbegin...\qfullend,
% so the WHOLE question (stem, figures, options) is set across the page in reading order, then the two columns
% resume. The question's head - its number, stem and first figure (\qfullsplit marks the end of the head) -
% is kept on one page, so a number is never left at a page foot with its figure overleaf; the rest may break
% across pages normally, so no big gaps are left. Without a \qfullsplit, the whole question is kept together
% if it is under 45% of a page. Never write these yourself.
\newif\ifqf@full
\newif\ifqf@boxed
\newsavebox{\qf@box}
\newcommand{\qf@open}{\global\qf@fulltrue\par\global\qf@boxedtrue
  \setbox\qf@box\vbox\bgroup\hsize\textwidth\linewidth\textwidth}
\newcommand{\qf@emit}[1]{\par\egroup\global\qf@boxedfalse
  \ifdim\dimexpr\ht\qf@box+\dp\qf@box\relax<#1\textheight\noindent\usebox{\qf@box}\par\else\unvbox\qf@box\fi}
\newcommand{\qfullsplit}{\ifqf@boxed\qf@emit{0.85}\fi}
\newcommand{\qf@close}{\ifqf@boxed\qf@emit{0.45}\fi\global\qf@fullfalse}
\newcommand{\qfullbegin}{\end{multicols}\qf@open}
\newcommand{\qfullend}{\qf@close\begin{multicols}{2}}
\newcommand{\qfullstart}{\qf@open}
\newcommand{\qfulldone}{\qf@close}
% \qfig[scale]{id}: block figure, centred, directly where it is written (size set by the tool).
% Fallback if a wide figure is not inside a full-width question: break out of the columns in place - no floats.
\newcommand{\qfig}[2][1]{\qf@check{#2}%
  \ifqf@full\qf@blk{#1}{#2}\else
  \ifnum\@nameuse{qf@wide@#2}=1 \expandafter\@firstoftwo\else\expandafter\@secondoftwo\fi
  {\end{multicols}\begin{center}\qf@img{#1}{\textwidth}{#2}\end{center}\begin{multicols}{2}}%
  {\qf@blk{#1}{#2}}\fi}
\newcommand{\qf@blk}[2]{\par\smallskip{\centering\qf@img{#1}{\linewidth}{#2}\par}\smallskip}
% \qfigi[scale]{id}: inline figure (inside option lists / tables / minipages); never breaks columns.
\newcommand{\qfigi}[2][1]{\qf@check{#2}\qf@img[valign=c]{#1}{\linewidth}{#2}}
\makeatother
% structure macros - lint reads these to check completeness and figure placement
\newcommand{\papersec}[1]{\par\bigskip{\centering\Large\bfseries PAPER-#1\par}\medskip}
\newcommand{\partsec}[2]{\par\medskip\noindent{\large\bfseries PART-#1: #2}\par\smallskip}
\newcommand{\sect}[1]{\par\smallskip\noindent\textbf{#1}\par}
\newcommand{\qn}[1]{\par\smallskip\noindent\textbf{#1.}~}
\newcommand{\paraq}[2]{\par\smallskip\noindent\textbf{Paragraph for Questions #1 and #2}\par}
\newenvironment{qoptions}{\begin{itemize}[leftmargin=1.8em,itemsep=1pt,topsep=2pt,label={}]}{\end{itemize}}
\newcommand{\opt}[2]{\item[(#1)] #2}
% figure registry (generated by score2_tools.py finalize)
\qfigdef{p1_phy_s1_q01_1}{Score Advanced - Test 41-img/p1_phy_s1_q01_1.png}{107.8}{0}
\qfigdef{p1_phy_s1_q03_1}{Score Advanced - Test 41-img/p1_phy_s1_q03_1.png}{88.0}{0}
\qfigdef{p1_phy_s1_q04_1}{Score Advanced - Test 41-img/p1_phy_s1_q04_1.png}{139.9}{0}
\qfigdef{p1_phy_s2_q05_1}{Score Advanced - Test 41-img/p1_phy_s2_q05_1.png}{95.0}{0}
\qfigdef{p1_phy_s2_q06_1}{Score Advanced - Test 41-img/p1_phy_s2_q06_1.png}{134.7}{0}
\qfigdef{p1_phy_s2_q07_1}{Score Advanced - Test 41-img/p1_phy_s2_q07_1.png}{161.3}{0}
\qfigdef{p1_phy_s3_q08_1}{Score Advanced - Test 41-img/p1_phy_s3_q08_1.png}{66.3}{0}
\qfigdef{p1_phy_s3_q10_1}{Score Advanced - Test 41-img/p1_phy_s3_q10_1.png}{71.3}{0}
\qfigdef{p1_phy_s3_q10_2}{Score Advanced - Test 41-img/p1_phy_s3_q10_2.png}{70.6}{0}
\qfigdef{p1_phy_s3_q10_3}{Score Advanced - Test 41-img/p1_phy_s3_q10_3.png}{71.3}{0}
\qfigdef{p1_phy_s3_q10_4}{Score Advanced - Test 41-img/p1_phy_s3_q10_4.png}{71.3}{0}
\qfigdef{p1_phy_s3_q10_5}{Score Advanced - Test 41-img/p1_phy_s3_q10_5.png}{71.3}{0}
\qfigdef{p1_phy_s3_q11_1}{Score Advanced - Test 41-img/p1_phy_s3_q11_1.png}{103.2}{0}
\qfigdef{p1_phy_s3_q11_2}{Score Advanced - Test 41-img/p1_phy_s3_q11_2.png}{104.4}{0}
\qfigdef{p1_phy_s3_q11_3}{Score Advanced - Test 41-img/p1_phy_s3_q11_3.png}{104.4}{0}
\qfigdef{p1_phy_s3_q11_4}{Score Advanced - Test 41-img/p1_phy_s3_q11_4.png}{104.4}{0}
\qfigdef{p1_phy_s3_q11_5}{Score Advanced - Test 41-img/p1_phy_s3_q11_5.png}{104.4}{0}
\qfigdef{p1_phy_s4_q01_1}{Score Advanced - Test 41-img/p1_phy_s4_q01_1.png}{150.4}{0}
\qfigdef{p1_phy_s4_q02_1}{Score Advanced - Test 41-img/p1_phy_s4_q02_1.png}{57.3}{0}
\qfigdef{p1_phy_s4_q03_1}{Score Advanced - Test 41-img/p1_phy_s4_q03_1.png}{161.1}{0}
\qfigdef{p1_phy_s4_q05_1}{Score Advanced - Test 41-img/p1_phy_s4_q05_1.png}{149.6}{0}
\qfigdef{p1_phy_s4_q06_1}{Score Advanced - Test 41-img/p1_phy_s4_q06_1.png}{65.1}{0}
\qfigdef{p1_chem_s1_q03_R}{Score Advanced - Test 41-img/p1_chem_s1_q03_R.png}{124.7}{0}
\qfigdef{p1_chem_s1_q03_A2}{Score Advanced - Test 41-img/p1_chem_s1_q03_A2.png}{64.1}{0}
\qfigdef{p1_chem_s1_q03_B2}{Score Advanced - Test 41-img/p1_chem_s1_q03_B2.png}{52.0}{0}
\qfigdef{p1_chem_s1_q03_C2}{Score Advanced - Test 41-img/p1_chem_s1_q03_C2.png}{62.9}{0}
\qfigdef{p1_chem_s1_q03_D2}{Score Advanced - Test 41-img/p1_chem_s1_q03_D2.png}{54.9}{0}
\qfigdef{p1_chem_s2_q06_1}{Score Advanced - Test 41-img/p1_chem_s2_q06_1.png}{133.2}{0}
\qfigdef{p1_chem_s2_q07_A2}{Score Advanced - Test 41-img/p1_chem_s2_q07_A2.png}{134.7}{0}
\qfigdef{p1_chem_s2_q07_B2}{Score Advanced - Test 41-img/p1_chem_s2_q07_B2.png}{130.6}{0}
\qfigdef{p1_chem_s2_q07_C2}{Score Advanced - Test 41-img/p1_chem_s2_q07_C2.png}{131.7}{0}
\qfigdef{p1_chem_s2_q07_D2}{Score Advanced - Test 41-img/p1_chem_s2_q07_D2.png}{121.0}{0}
\qfigdef{p1_chem_s3_q11_X}{Score Advanced - Test 41-img/p1_chem_s3_q11_X.png}{62.4}{0}
\qfigdef{p1_chem_s3_q11_P}{Score Advanced - Test 41-img/p1_chem_s3_q11_P.png}{100.7}{0}
\qfigdef{p1_chem_s3_q11_Q}{Score Advanced - Test 41-img/p1_chem_s3_q11_Q.png}{100.7}{0}
\qfigdef{p1_chem_s3_q11_R}{Score Advanced - Test 41-img/p1_chem_s3_q11_R.png}{100.7}{0}
\qfigdef{p1_chem_s3_q11_S}{Score Advanced - Test 41-img/p1_chem_s3_q11_S.png}{100.7}{0}
\qfigdef{p1_chem_s4_q05_1}{Score Advanced - Test 41-img/p1_chem_s4_q05_1.png}{151.4}{0}
\qfigdef{p1_chem_s4_q06_1}{Score Advanced - Test 41-img/p1_chem_s4_q06_1.png}{245.2}{0}
\qfigdef{p2_phy_s1_q01_1}{Score Advanced - Test 41-img/p2_phy_s1_q01_1.png}{50.4}{0}
\qfigdef{p2_phy_s1_q03_A}{Score Advanced - Test 41-img/p2_phy_s1_q03_A.png}{106.5}{0}
\qfigdef{p2_phy_s1_q03_B}{Score Advanced - Test 41-img/p2_phy_s1_q03_B.png}{106.6}{0}
\qfigdef{p2_phy_s1_q03_C}{Score Advanced - Test 41-img/p2_phy_s1_q03_C.png}{106.6}{0}
\qfigdef{p2_phy_s1_q03_D}{Score Advanced - Test 41-img/p2_phy_s1_q03_D.png}{106.5}{0}
\qfigdef{p2_phy_s1_q04_1}{Score Advanced - Test 41-img/p2_phy_s1_q04_1.png}{84.6}{0}
\qfigdef{p2_phy_s2_q05_1}{Score Advanced - Test 41-img/p2_phy_s2_q05_1.png}{81.8}{0}
\qfigdef{p2_phy_s2_q05_A}{Score Advanced - Test 41-img/p2_phy_s2_q05_A.png}{90.1}{0}
\qfigdef{p2_phy_s2_q05_B}{Score Advanced - Test 41-img/p2_phy_s2_q05_B.png}{87.2}{0}
\qfigdef{p2_phy_s2_q05_C}{Score Advanced - Test 41-img/p2_phy_s2_q05_C.png}{87.0}{0}
\qfigdef{p2_phy_s2_q05_D}{Score Advanced - Test 41-img/p2_phy_s2_q05_D.png}{87.2}{0}
\qfigdef{p2_phy_s2_q07_1}{Score Advanced - Test 41-img/p2_phy_s2_q07_1.png}{70.5}{0}
\qfigdef{p2_phy_s2_q07_A}{Score Advanced - Test 41-img/p2_phy_s2_q07_A.png}{57.5}{0}
\qfigdef{p2_phy_s2_q07_B}{Score Advanced - Test 41-img/p2_phy_s2_q07_B.png}{59.6}{0}
\qfigdef{p2_phy_s2_q07_C}{Score Advanced - Test 41-img/p2_phy_s2_q07_C.png}{58.0}{0}
\qfigdef{p2_phy_s2_q07_D}{Score Advanced - Test 41-img/p2_phy_s2_q07_D.png}{57.6}{0}
\qfigdef{p2_phy_s3_para1_2_1}{Score Advanced - Test 41-img/p2_phy_s3_para1_2_1.png}{108.5}{0}
\qfigdef{p2_phy_s3_para3_4_1}{Score Advanced - Test 41-img/p2_phy_s3_para3_4_1.png}{123.9}{0}
\qfigdef{p2_phy_s4_q05_1}{Score Advanced - Test 41-img/p2_phy_s4_q05_1.png}{113.7}{0}
\qfigdef{p2_phy_s4_q06_1}{Score Advanced - Test 41-img/p2_phy_s4_q06_1.png}{68.3}{0}
\qfigdef{p2_chem_s1_q03_R}{Score Advanced - Test 41-img/p2_chem_s1_q03_R.png}{118.0}{0}
\qfigdef{p2_chem_s1_q03_A2}{Score Advanced - Test 41-img/p2_chem_s1_q03_A2.png}{21.0}{0}
\qfigdef{p2_chem_s1_q03_B2}{Score Advanced - Test 41-img/p2_chem_s1_q03_B2.png}{34.8}{0}
\qfigdef{p2_chem_s1_q03_C2}{Score Advanced - Test 41-img/p2_chem_s1_q03_C2.png}{46.0}{0}
\qfigdef{p2_chem_s1_q03_D2}{Score Advanced - Test 41-img/p2_chem_s1_q03_D2.png}{38.5}{0}
\qfigdef{p2_chem_s1_q04_R}{Score Advanced - Test 41-img/p2_chem_s1_q04_R.png}{126.9}{0}
\qfigdef{p2_chem_s1_q04_B2}{Score Advanced - Test 41-img/p2_chem_s1_q04_B2.png}{31.8}{0}
\qfigdef{p2_chem_s1_q04_D2}{Score Advanced - Test 41-img/p2_chem_s1_q04_D2.png}{34.5}{0}
\qfigdef{p2_chem_s2_q06_X}{Score Advanced - Test 41-img/p2_chem_s2_q06_X.png}{52.0}{0}
\qfigdef{p2_chem_s4_q08_T}{Score Advanced - Test 41-img/p2_chem_s4_q08_T.png}{37.2}{0}
\qfigdef{p2_chem_s4_q08_i}{Score Advanced - Test 41-img/p2_chem_s4_q08_i.png}{107.2}{0}
\qfigdef{p2_chem_s4_q08_ii}{Score Advanced - Test 41-img/p2_chem_s4_q08_ii.png}{104.4}{0}
\qfigdef{p2_chem_s4_q08_iii}{Score Advanced - Test 41-img/p2_chem_s4_q08_iii.png}{106.1}{0}
\qfigdef{p2_chem_s4_q08_iv}{Score Advanced - Test 41-img/p2_chem_s4_q08_iv.png}{112.3}{0}
\qfigdef{p2_chem_s4_q08_v}{Score Advanced - Test 41-img/p2_chem_s4_q08_v.png}{85.4}{0}
\qfigdef{p2_chem_s4_q09_i}{Score Advanced - Test 41-img/p2_chem_s4_q09_i.png}{150.3}{0}
\qfigdef{p2_chem_s4_q09_ii}{Score Advanced - Test 41-img/p2_chem_s4_q09_ii.png}{148.6}{0}
\qfigdef{p2_chem_s4_q09_iii}{Score Advanced - Test 41-img/p2_chem_s4_q09_iii.png}{110.2}{0}
\qfigdef{p2_chem_s4_q09_iv}{Score Advanced - Test 41-img/p2_chem_s4_q09_iv.png}{140.1}{0}
\qfigdef{p2_chem_s4_q09_v}{Score Advanced - Test 41-img/p2_chem_s4_q09_v.png}{120.0}{0}
\qfigdef{p2_chem_s4_q10_i}{Score Advanced - Test 41-img/p2_chem_s4_q10_i.png}{76.9}{0}
\qfigdef{p2_chem_s4_q10_ii}{Score Advanced - Test 41-img/p2_chem_s4_q10_ii.png}{68.5}{0}
\qfigdef{p2_chem_s4_q10_iii}{Score Advanced - Test 41-img/p2_chem_s4_q10_iii.png}{105.8}{0}
\qfigdef{p2_chem_s4_q10_iv}{Score Advanced - Test 41-img/p2_chem_s4_q10_iv.png}{98.4}{0}
\qfigdef{p2_chem_s4_q10_v}{Score Advanced - Test 41-img/p2_chem_s4_q10_v.png}{88.6}{0}
\qfigdef{p2_chem_s4_q10_vi}{Score Advanced - Test 41-img/p2_chem_s4_q10_vi.png}{94.3}{0}
\raggedcolumns
\makeatletter\@beginparpenalty=10000\makeatother
\newcolumntype{L}[1]{>{\raggedright\arraybackslash}p{#1}}
\begin{document}
\begin{center}{\Large\bfseries Score Advanced --- Test 41}\end{center}
\begin{multicols}{2}
\raggedcolumns
\papersec{1}
\partsec{1}{PHYSICS}
\sect{SECTION-I (i)}
\qn{1} Figure shows a metal ball suspended by thread of negligible mass from an upright cylinder that floats partially submerged in water. The cylinder has height 6 cm, face area 11 cm$^2$ on the top and bottom and density 0.5 g/cm$^3$. 4 cm of cylinder's height is inside the water surface. If density of the metal ball is 8 gm/cm$^3$ then its radius is equal to :- ($\rho_\mathrm w$ = 1 gm/cm$^3$)\par
\qfig{p1_phy_s1_q01_1}
\begin{qoptions}
\opt{A}{$\left(\dfrac38\right)^{1/3}$ cm}
\opt{B}{$\left(\dfrac34\right)^{1/3}$ cm}
\opt{C}{$\left(\dfrac43\right)^{1/3}$ cm}
\opt{D}{$\left(\dfrac5{11}\right)^{1/3}$ cm}
\end{qoptions}
\qn{2} A light ray is incident on a transparent sphere of index = $\sqrt2$, at an angle of incidence = 45$^\circ$. What is the deviation of a tiny fraction of the ray, which enters the sphere, undergoes two internal reflections and then refracts out into air ?
\begin{qoptions}
\opt{A}{270$^\circ$}
\opt{B}{240$^\circ$}
\opt{C}{120$^\circ$}
\opt{D}{180$^\circ$}
\end{qoptions}
\qn{3} In an aluminum (Al) bar of square cross section, a square hole is drilled and is filled with iron (Fe) as shown in the figure. The electrical resistivities of Al and Fe are 2.7$\times$10$^{-8}\,\Omega$m and 1.0 $\times$ 10$^{-7}$ $\Omega$m, respectively. The electrical resistance between the two faces P and Q of the composite bar is :-\par
\qfig{p1_phy_s1_q03_1}
\begin{qoptions}
\opt{A}{$\dfrac{2475}{64}\,\mu\Omega$}
\opt{B}{$\dfrac{1875}{64}\,\mu\Omega$}
\opt{C}{$\dfrac{1875}{49}\,\mu\Omega$}
\opt{D}{$\dfrac{2475}{132}\,\mu\Omega$}
\end{qoptions}
\qn{4} In the figure shown AB is a rod of length 30 cm and area of cross-section 1.0 cm$^2$ and thermal conductivity 336 S. I. units. The ends A \& B are maintained at temperatures 20$^\circ$C and 40$^\circ$C respectively. A point C of this rod is connected to a box D, containing ice at 0$^\circ$C, through a highly conducting wire of negligible heat capacity. The rate at which ice melts in the box is :-\par
[Assume latent heat of fusion for ice L = 80 cal/gm]\par
\qfig{p1_phy_s1_q04_1}
\begin{qoptions}
\opt{A}{84 mg/s}
\opt{B}{84 g/s}
\opt{C}{20 mg/s}
\opt{D}{40 mg/s}
\end{qoptions}
\sect{SECTION-I (ii)}
\qn{5} A thermally insulated chamber of volume 2V$_0$ is divided by a frictionless piston of area S into two equal parts A and B. Part A has an ideal gas at pressure P$_0$ and temperature T$_0$ and in part B is vacuum. A massless spring of force constant k is connected with a massless piston and the wall of the container as shown. Initially spring is unstretched. Gas in chamber A is allowed to expand. Let in equilibrium spring is compressed by x$_0$. Then :-\par
\qfig{p1_phy_s2_q05_1}
\begin{qoptions}
\opt{A}{Pressure of the gas at equilibrium is $\dfrac{\mathrm{kx}_0}{\mathrm S}$.}
\opt{B}{Work done by the gas upto equilibrium is $\dfrac12$ kx$^2_0$.}
\opt{C}{Change in internal energy of the gas upto equilibrium is $\dfrac12$ kx$^2_0$.}
\opt{D}{At equilibrium temperature of the gas is decreased.}
\end{qoptions}
\qn{6} A uniform magnetic field B exists in the region between x = 0 and x = $\dfrac{3\mathrm R}2$ (region 2 in the figure) pointing normally into the plane of the paper. A particle with charge +Q and momentum p directed along x-axis enters region 2 from region 1 at point P$_1$ (y = $-$R). Which of the following option(s) is/are correct?\par
\qfig{p1_phy_s2_q06_1}
\begin{qoptions}
\opt{A}{For B = $\dfrac8{13}\dfrac{\mathrm p}{\mathrm{QR}}$, the particle will enter region 3 through the point P$_2$ on x-axis.}
\opt{B}{For B $>$ $\dfrac23\dfrac{\mathrm p}{\mathrm{QR}}$, the particle will re-enter region 1.}
\opt{C}{For a fixed B, particles of same charge Q and same velocity v, the distance between the point P$_1$ and the point of re-entry into region 1 is inversely proportional to the mass of the particle.}
\opt{D}{When the particle re-enters region 1 through the longest possible path in region 2, the magnitude of the change in its linear momentum between point P$_1$ and the farthest point from y-axis is p/$\sqrt2$.}
\end{qoptions}
\qn{7} Two blocks of masses 3 kg and 6 kg rest on a horizontal smooth surface. The 3 kg block is attached to a spring with a force constant k = 900 Nm$^{-1}$ which is compressed 2 m from the equilibrium position. The 6 kg block is at rest at 1 m from mean position. 3 kg mass strikes the 6 kg mass and the two stick together.\par
\qfig{p1_phy_s2_q07_1}
\begin{qoptions}
\opt{A}{Velocity of the combined masses immediately after the collision is 10 ms$^{-1}$.}
\opt{B}{Velocity of the combined masses immediately after the collision is 5 ms$^{-1}$.}
\opt{C}{Amplitude of the resulting oscillation is $\sqrt2$ m.}
\opt{D}{Amplitude of the resulting oscillation is $\sqrt5$/2m.}
\end{qoptions}
\sect{SECTION-I (iii)}
\qn{8} A block of mass m is stationary with respect to a rough wedge as shown in figure. Starting from rest, in time t work done on the block :  (m = 1kg, $\theta$ = 30$^\circ$, a = 2m/s$^2$, t = 4s)\par
\qfig{p1_phy_s3_q08_1}
{\footnotesize\setlength{\tabcolsep}{1.5pt}\noindent\begin{tabular}{|c|L{3.2cm}|c|L{1.6cm}|}\hline & \multicolumn{1}{>{\centering\arraybackslash}p{3.2cm}|}{List-I} & & \multicolumn{1}{>{\centering\arraybackslash}p{1.6cm}|}{List-II} \\ \hline
(P) & By gravity & (1) & 144 J \\ \hline
(Q) & By normal reaction & (2) & 32 J \\ \hline
(R) & By friction & (3) & 56 J \\ \hline
(S) & By all the forces & (4) & 48 J \\ \hline
 & & (5) & None \\ \hline
\end{tabular}}\par
\begin{qoptions}
\opt{A}{P $\to$ 5;Q $\to$ 1;R $\to$ 4;S $\to$ 2}
\opt{B}{P $\to$ 4;Q $\to$ 3;R $\to$ 2;S $\to$ 1}
\opt{C}{P $\to$ 1;Q $\to$ 2;R $\to$ 3;S $\to$ 5}
\opt{D}{P $\to$ 4;Q $\to$ 2;R $\to$ 3;S $\to$ 1}
\end{qoptions}
\qn{9} An isolated zinc bead of radius r charged to a negative potential V$_0$ (assuming potential at infinitely distant points to be zero) placed in free space is being continuously irradiated by ultraviolet light of wavelength $\lambda$. Photoelectric threshold wavelength for zinc is $\lambda_0$, speed of light is c, Denoting Planck's constant by h, permittivity of free space by $\varepsilon_0$, charge and mass of an electron by $-$e and m.\par
{\footnotesize\setlength{\tabcolsep}{1.5pt}\noindent\begin{tabular}{|c|L{2.0cm}|c|L{3.4cm}|}\hline & \multicolumn{1}{>{\centering\arraybackslash}p{2.0cm}|}{List-I} & & \multicolumn{1}{>{\centering\arraybackslash}p{3.4cm}|}{List-II} \\ \hline
(P) & Maximum speed of a photoelectron immediately after emission is & (1) & $\dfrac{\mathrm{hc}}{\mathrm e}\left(\dfrac1\lambda-\dfrac1{\lambda_0}\right)$ \\ \hline
(Q) & Maximum speed of a photoelectron at a great distance from the bead is & (2) & $\sqrt{\dfrac{2\mathrm{hc}}{\mathrm m}\left(\dfrac1\lambda-\dfrac1{\lambda_0}\right)+\dfrac{\mathrm{eV}_0}{\mathrm m}}$ \\ \hline
(R) & Potential of the bead after prolonged irradiation is & (3) & $\dfrac{4\pi\varepsilon_0\mathrm r}{\mathrm e}\Bigl\{\dfrac{\mathrm{hc}}{\mathrm e}\Bigl(\dfrac1\lambda-\dfrac1{\lambda_0}\Bigr)+\mathrm V_0\Bigr\}$ \\ \hline
(S) & Total number of photoelectrons emitted from the bead is & (4) & $\sqrt{\dfrac{2\mathrm{hc}}{\mathrm m}\left(\dfrac1\lambda-\dfrac1{\lambda_0}\right)}$ \\ \hline
 & & (5) & None of these \\ \hline
\end{tabular}}\par
\begin{qoptions}
\opt{A}{P $\to$ 3;Q $\to$ 2;R $\to$ 1;S $\to$ 4}
\opt{B}{P $\to$ 4;Q $\to$ 5;R $\to$ 1;S $\to$ 3}
\opt{C}{P $\to$ 5;Q $\to$ 2;R $\to$ 1;S $\to$ 4}
\opt{D}{P $\to$ 3;Q $\to$ 2;R $\to$ 4;S $\to$ 1}
\end{qoptions}
\qn{10} Two transparent media of refractive indices $\mu_1$ and $\mu_3$ have a solid lens shaped transparent material of refractive index $\mu_2$ between them as shown in figures in List$-$II. A ray traversing these media is also shown in the figures. In List$-$I different relationships between $\mu_1$, $\mu_2$ and $\mu_3$ are given. Match them to the ray diagrams shown in List$-$II.\par
{\footnotesize\setlength{\tabcolsep}{1.5pt}\noindent\begin{tabular}{|c|L{1.6cm}|c|c|}\hline & \multicolumn{1}{>{\centering\arraybackslash}p{1.6cm}|}{List-I} & & \multicolumn{1}{c|}{List-II} \\ \hline
(P) & $\mu_1<\mu_2$ & (1) & \qfigi{p1_phy_s3_q10_1} \\ \hline
(Q) & $\mu_1>\mu_2$ & (2) & \qfigi{p1_phy_s3_q10_2} \\ \hline
(R) & $\mu_2=\mu_3$ & (3) & \qfigi{p1_phy_s3_q10_3} \\ \hline
(S) & $\mu_2>\mu_3$ & (4) & \qfigi{p1_phy_s3_q10_4} \\ \hline
 & & (5) & \qfigi{p1_phy_s3_q10_5} \\ \hline
\end{tabular}}\par
\begin{qoptions}
\opt{A}{P $\to$ 1,4;Q $\to$ 2,3;R $\to$ 3,4;S $\to$ 1,5}
\opt{B}{P $\to$ 1,4,5;Q $\to$ 1,3,5;R $\to$ 1,2;S $\to$ 1,5}
\opt{C}{P $\to$ 1,3;Q $\to$ 2,4,5;R $\to$ 1,3,5;S $\to$ 2,4}
\opt{D}{P $\to$ 1,3,5;Q $\to$ 2,4;R $\to$ 2,3;S $\to$ 4,5}
\end{qoptions}
\qn{11} You are given many resistances, capacitors and inductors. These are connected to a variable DC voltage source (the first two circuits) or an AC voltage source of 50 Hz frequency (the next three circuits) in different ways as shown in List$-$II. When a current I (steady state for DC or rms for AC) flows through the circuit, the corresponding voltage V$_1$ and V$_2$ (indicated in circuits) are related as shown in List$-$I. Match the two list.\par
{\footnotesize\setlength{\tabcolsep}{1.5pt}\noindent\begin{tabular}{|c|L{2.0cm}|c|c|}\hline & \multicolumn{1}{>{\centering\arraybackslash}p{2.0cm}|}{List-I} & & \multicolumn{1}{c|}{List-II} \\ \hline
(P) & I $\neq$ 0, V$_1$ is proportional to I & (1) & \qfigi{p1_phy_s3_q11_1} \\ \hline
(Q) & I $\neq$ 0, V$_2>$ V$_1$ & (2) & \qfigi{p1_phy_s3_q11_2} \\ \hline
(R) & V$_1$=0, V$_2$=V & (3) & \qfigi{p1_phy_s3_q11_3} \\ \hline
(S) & I $\neq$ 0, V$_2$ is proportional to I & (4) & \qfigi{p1_phy_s3_q11_4} \\ \hline
 & & (5) & \qfigi{p1_phy_s3_q11_5} \\ \hline
\end{tabular}}\par
\begin{qoptions}
\opt{A}{P $\to$ 3,4,5;Q $\to$ 2,3,4,5;R $\to$ 1,2;S $\to$ 2,3,4,5}
\opt{B}{P $\to$ 1,3,5;Q $\to$ 2,3,4,5;R $\to$ 3,4;S $\to$ 2,5}
\opt{C}{P $\to$ 3,4,5;Q $\to$ 2;R $\to$ 3,5;S $\to$ 1,4}
\opt{D}{P $\to$ 2,4;Q $\to$ 2,3;R $\to$ 1,4;S $\to$ 1,5}
\end{qoptions}
\sect{SECTION-II}
\qn{1} A 2 $\mu$F capacitor C$_1$ is first charged to a potential difference of 10 V using a battery. Then the battery is removed and the capacitor is connected to an uncharged capacitor C$_2$ of 8$\mu$F. The charge in C$_2$ on equilibrium condition is\underline{\hspace{1cm}}$\mu$C. (Round off to the Nearest Integer)\par
\qfig{p1_phy_s4_q01_1}
\qn{2} A bullet of 10 g, moving with velocity v, collides head-on with the stationary bob of a pendulum and recoils with velocity 100 m/s. The length of the pendulum is 0.5 m and mass of the bob is 1 kg. The minimum value of v = \underline{\hspace{1cm}} m/s so that the pendulum describes a circle. (Assume the string to be inextensible and g = 10 m/s$^2$)\par
\qfig{p1_phy_s4_q02_1}
\qn{3} An electric field $\vec{\mathrm E}$ = 4x$\hat{\mathrm i}$ $-$ (y$^2$ + 1)$\hat{\mathrm j}$ N/C passes through the box shown in figure. The flux of the electric field through surfaces ABCD and BCGF are marked as $\phi_\mathrm I$ and $\phi_\mathrm{II}$ respectively. The magnitude of difference between ($\phi_\mathrm I$ $-$ $\phi_\mathrm{II}$) is (in Nm$^2$/C) \underline{\hspace{1.5cm}}.\par
\qfig{p1_phy_s4_q03_1}
\qn{4} A signal of 0.1 kW is transmitted in a cable. The attenuation of cable is $-$5 dB per km and cable length is 20 km. The power received at receiver is 10$^{-x}$ W. The value of x is \underline{\hspace{1cm}}.\par
[Gain in dB = 10 log$_{10}\left(\dfrac{\mathrm P_0}{\mathrm P_\mathrm i}\right)$]
\qn{5} A carpet of mass 'M' made of inextensible material is rolled along its length in the form of a cylinder of radius 'R' and is kept on a rough floor. The carpet starts unrolling without sliding on the floor when a negligibly small push is given to it. The horizontal velocity of the axis of the cylindrical part of the carpet when its radius reduces to R/2 is $\sqrt{\dfrac{\mathrm x\times7\mathrm{gR}}3}$ then find x.\par
\qfig{p1_phy_s4_q05_1}
\qn{6} A particle is projected from point A, that is at a distance 4R from the centre of the Earth, with speed v$_1$ in a direction making 30$^\circ$ with the line joining the centre of the Earth and point A, as shown. Find the speed v$_1$ if particle passes grazing the surface of the earth. Consider gravitational interaction only between these two. (use $\dfrac{\mathrm{GM}}{\mathrm R}$ = 6.4 $\times$ 10$^7$ m$^2$/s$^2$) Express your answer in the form $\left(500\sqrt2\,\mathrm X\right)$ m/s and write the value of X.\par
\qfig{p1_phy_s4_q06_1}
\raggedcolumns
\papersec{1}
\partsec{2}{CHEMISTRY}
\sect{SECTION-I (i)}
\qn{1} Identify the incorrect statement(s) among the following :\par
(i) Moist ammonia gas can be dried by using anhydrous CaCl$_2$.\par
(ii) When metallic copper react with very dilute nitric acid, ammonium nitrite is one of the products.\par
(iii) Phosphine is a weaker base than ammonia because phosphorous atom is less electronegative than nitrogen.\par
(iv) NaH$_2$PO$_2$ is an acidic salt and it is named as sodium hypophosphate.
\begin{qoptions}
\opt{A}{(i) and (iii)}
\opt{B}{(i) and (iv)}
\opt{C}{(ii) and (iii)}
\opt{D}{All of these}
\end{qoptions}
\qn{2} 2.5 mL of $\dfrac25$ M weak monoacidic base (K$_\mathrm b$ = 1 $\times$ 10$^{-12}$ at 25$^\circ$C) is titrated with $\dfrac2{15}$ M HCl in water at 25$^\circ$C. The concentration of H$^+$ at equivalence point is\par
(K$_\mathrm w$ = 1 $\times$ 10$^{-14}$ at 25$^\circ$C)
\begin{qoptions}
\opt{A}{3.7 $\times$ 10$^{-13}$ M}
\opt{B}{3.2 $\times$ 10$^{-7}$ M}
\opt{C}{3.2 $\times$ 10$^{-2}$ M}
\opt{D}{2.7 $\times$ 10$^{-2}$ M}
\end{qoptions}
\qn{3} Sodium phenoxide when heated with CO$_2$ under pressure at 125$^\circ$C yields a product which on acetylation produces C.\par
\qfig{p1_chem_s1_q03_R}
The major product C would be :
\begin{qoptions}
\opt{A}{\qfigi{p1_chem_s1_q03_A2}}
\opt{B}{\qfigi{p1_chem_s1_q03_B2}}
\opt{C}{\qfigi{p1_chem_s1_q03_C2}}
\opt{D}{\qfigi{p1_chem_s1_q03_D2}}
\end{qoptions}
\qn{4} Nitric acid is prepared by ostwald method using given Steps : -\par
NH$_3$+ O$_2$ $\xrightarrow{\text{catalyst}}$ (X) + H$_2$O\par
(X) + O$_2$ $\to$ (Y)\par
(Y) + H$_2$O $\to$ HNO$_3$ + (X)\par
(X) and (Y) are respectively
\begin{qoptions}
\opt{A}{NO$_2$ and NO}
\opt{B}{NO and N$_2$O}
\opt{C}{N$_2$O and NO$_2$}
\opt{D}{NO and NO$_2$}
\end{qoptions}
\sect{SECTION-I (ii)}
\qn{5} The number of correct statements among the following is
\begin{qoptions}
\opt{A}{NF$_3$ $<$ NH$_3$ $<$ CH$_4$ $<$ CO$_2$ (bond angle)}
\opt{B}{S$-$H----O $<$ O$-$H----S $<$ F$-$H----O $<$ F$-$H----F$^{-}$ (Strength of H-bond)}
\opt{C}{Be$_2$H$_4$ = C$_2$H$_4$ $<$ B$_2$H$_6$ (max number of H atoms in plane)}
\opt{D}{O$_2$ $<$ O$_2^-$ $<$ O$_2^{-2}$ (unpaired electron)}
\end{qoptions}
\qn{6} \qfigi{p1_chem_s2_q06_1}\par
G \& reagent can be :
\begin{qoptions}
\opt{A}{$-$CN ; SnCl$_2$ + HCl (aq.)}
\opt{B}{$-$COCl ; H$_2$ (Pd / BaSO$_4$)}
\opt{C}{$-$COOEt ; DIBALH ($-$78$^\circ$C)}
\opt{D}{$-$CH = CH$_2$ ; O$_3$ / Zn / H$_2$O}
\end{qoptions}
\qn{7} Which of the following is/are reducing sugar/s ?
\begin{qoptions}
\opt{A}{\qfigi{p1_chem_s2_q07_A2}}
\opt{B}{\qfigi{p1_chem_s2_q07_B2}}
\opt{C}{\qfigi{p1_chem_s2_q07_C2}}
\opt{D}{\qfigi{p1_chem_s2_q07_D2}}
\end{qoptions}
\sect{SECTION-I (iii)}
\qn{8} Match the reactions (in the given stoichiometry of the reactants) in List-I with one of their products given in List-II and choose the correct option.\par
{\footnotesize\setlength{\tabcolsep}{1.5pt}\noindent\begin{tabular}{|c|L{3.2cm}|c|L{2.2cm}|}\hline & \multicolumn{1}{>{\centering\arraybackslash}p{3.2cm}|}{List-I} & & \multicolumn{1}{>{\centering\arraybackslash}p{2.2cm}|}{List-II} \\ \hline
(P) & P$_2$O$_3$ + 3H$_2$O $\longrightarrow$ & (1) & P(O)(OCH$_3$)Cl$_2$ \\ \hline
(Q) & P$_4$ + 3NaOH + 3H$_2$O $\longrightarrow$ & (2) & H$_3$PO$_3$ \\ \hline
(R) & PCl$_5$ + CH$_3$COOH $\longrightarrow$ & (3) & PH$_3$ \\ \hline
(S) & H$_3$PO$_2$ + 2H$_2$O + 4AgNO$_3$ $\longrightarrow$ & (4) & POCl$_3$ \\ \hline
 & & (5) & H$_3$PO$_4$ \\ \hline
\end{tabular}}\par
\begin{qoptions}
\opt{A}{P $\to$ 2;Q $\to$ 3;R $\to$ 1;S $\to$ 5}
\opt{B}{P $\to$ 3;Q $\to$ 5;R $\to$ 4;S $\to$ 2}
\opt{C}{P $\to$ 5;Q $\to$ 2;R $\to$ 1;S $\to$ 3}
\opt{D}{P $\to$ 2;Q $\to$ 3;R $\to$ 4;S $\to$ 5}
\end{qoptions}
\qn{9} Match the List :\par
{\footnotesize\setlength{\tabcolsep}{1.5pt}\noindent\begin{tabular}{|c|L{3.4cm}|c|L{2.6cm}|}\hline & \multicolumn{1}{>{\centering\arraybackslash}p{3.4cm}|}{List-I} & & \multicolumn{1}{>{\centering\arraybackslash}p{2.6cm}|}{List-II} \\ \hline
(P) & Tetragonal and Hexagonal & (1) & are two crystal systems \\ \hline
(Q) & Cubic and Rhombohedral & (2) & a = b $\neq$ c \\ \hline
(R) & Monoclinic and Triclinic & (3) & a $\neq$ b $\neq$ c \\ \hline
(S) & Cubic and Hexagonal & (4) & a = b = c \\ \hline
 & & (5) & $\alpha$ = $\beta$ = 90$^\circ$ \\ \hline
\end{tabular}}\par
\begin{qoptions}
\opt{A}{P $\to$ 1,2;Q $\to$ 1,2,5;R $\to$ 1,4;S $\to$ 1,5}
\opt{B}{P $\to$ 1,5;Q $\to$ 1,3;R $\to$ 1,4;S $\to$ 2,5}
\opt{C}{P $\to$ 1,3,5;Q $\to$ 1,2;R $\to$ 1,3;S $\to$ 1,2,5}
\opt{D}{P $\to$ 1,2,5;Q $\to$ 1,4;R $\to$ 1,3;S $\to$ 1,5}
\end{qoptions}
\qn{10} t$_\mathrm{x/y}$ = time in which x/y fraction of reactant converts into product :\par
{\footnotesize\setlength{\tabcolsep}{1.5pt}\noindent\begin{tabular}{|c|L{0.9cm}|c|L{4.9cm}|}\hline & \multicolumn{1}{>{\centering\arraybackslash}p{0.9cm}|}{List-I} & & \multicolumn{1}{>{\centering\arraybackslash}p{4.9cm}|}{List-II} \\ \hline
(P) & t$_{5/9}$ & (1) & Equal to 54 sec. if t$_{1/3}$ is 18 sec. in case of first order reaction. \\ \hline
(Q) & t$_{19/27}$ & (2) & Equal to 32 sec. if t$_{1/4}$ is 16 sec. in case of first order reaction. \\ \hline
(R) & t$_{7/8}$ & (3) & Equal to 56 sec. if t$_{1/3}$ is 4 sec. in case of second order reaction. \\ \hline
(S) & t$_{7/16}$ & (4) & Equal to 30 sec. if t$_{1/3}$ is 18 sec. in case of zero order reaction. \\ \hline
 & & (5) & Equal to 28 sec. if t$_{1/2}$ is 16 sec. in case of zero order reaction. \\ \hline
\end{tabular}}\par
\begin{qoptions}
\opt{A}{P $\to$ 1;Q $\to$ 2;R $\to$ 3;S $\to$ 4,5}
\opt{B}{P $\to$ 4;Q $\to$ 1;R $\to$ 3,5;S $\to$ 2}
\opt{C}{P $\to$ 2;Q $\to$ 3;R $\to$ 1,3;S $\to$ 5}
\opt{D}{P $\to$ 1,3;Q $\to$ 2,4;R $\to$ 3;S $\to$ 5}
\end{qoptions}
\qn{11} The desired product X can be prepared by reacting the major product of the reactions in List-I with one or more appropriate reagents in List-II.\par
(given, order of migratory aptitude: aryl $>$ alkyl $>$ hydrogen)\par
\qfig{p1_chem_s3_q11_X}
{\footnotesize\setlength{\tabcolsep}{1.5pt}\noindent\begin{tabular}{|c|L{3.3cm}|c|L{2.6cm}|}\hline & \multicolumn{1}{>{\centering\arraybackslash}p{3.3cm}|}{List-I} & & \multicolumn{1}{>{\centering\arraybackslash}p{2.6cm}|}{List-II} \\ \hline
(P) & \qfigi{p1_chem_s3_q11_P} & (1) & I$_2$, NaOH \\ \hline
(Q) & \qfigi{p1_chem_s3_q11_Q} & (2) & [Ag(NH$_3$)$_2$]OH \\ \hline
(R) & \qfigi{p1_chem_s3_q11_R} & (3) & Fehling solution \\ \hline
(S) & \qfigi{p1_chem_s3_q11_S} & (4) & HCHO, NaOH \\ \hline
 & & (5) & NaOBr \\ \hline
\end{tabular}}\par
The correct option is
\begin{qoptions}
\opt{A}{P $\to$ 1;Q $\to$ 2,3;R $\to$ 1,4;S $\to$ 2,4}
\opt{B}{P $\to$ 1,5;Q $\to$ 3,4;R $\to$ 4,5;S $\to$ 3}
\opt{C}{P $\to$ 1,5;Q $\to$ 1,5;R $\to$ 1,5;S $\to$ 1,5}
\opt{D}{P $\to$ 1,5;Q $\to$ 2,3;R $\to$ 1,5;S $\to$ 2,3}
\end{qoptions}
\sect{SECTION-II}
\qn{1} How many of the following have d$^2$sp$^3$ hybridization?\par
[Ni(NH$_3$)$_6$]$^{2+}$, [Co(NH$_3$)$_6$]$^{3+}$, [Fe(CN)$_6$]$^{4-}$,\par
[Cr(NH$_3$)$_6$]$^{3+}$, [Co(en)$_3$]$^{3+}$, [CoF$_6$]$^{3-}$, [Co(H$_2$O)$_6$]$^{2+}$
\qn{2} K$_2$Cr$_2$O$_7$ (Orange solid) $\xrightarrow{\Delta}$ ``X'' (Yellow compound) + `Y' (Green compound) + O$_2\uparrow$\par
In the given reaction how many statement is CORRECT?\par
(A) `Y' is CrO$_3$\par
(B) Aqueous solution of `X' with conc. H$_2$SO$_4$ is converted into orange solution\par
(C) When aqueous solution of `X' is treated with H$_2$O$_2$ in acidic medium blue colour compound is obtained\par
(D) `Y' is formed on thermal decomposition of (NH$_4$)$_2$Cr$_2$O$_7$
\qn{3} Consider the following reversible reaction,\par
A(g) + B(g) $\rightleftharpoons$ AB(g)\par
The activation energy of the backward reaction exceeds that of the forward reaction by 2RT (in J mol$^{-1}$). If the pre-exponential factor of the forward reaction is 4 times that of the reverse reaction, the absolute value of $\Delta$G$^\theta$ (in J mol$^{-1}$) for the reaction at 300 K is .........\par
(Given : ln (2) = 0.7, RT = 2500 J mol$^{-1}$ at 300 K and G is the Gibbs energy)
\qn{4} Estimate the cell potential (in milli volt) of a Daniel cell having 1M Zn$^{++}$ and originally having 1M Cu$^{++}$ after sufficient NH$_3$ has been added to the cathode compartment to make NH$_3$ concentration 2M. K$_\mathrm f$ for [Cu(NH$_3$)$_4$]$^{2+}$ = 1 $\times$ 10$^{12}$, E$^\circ$ for the reaction,\par
Zn + Cu$^{2+}$ $\longrightarrow$ Zn$^{2+}$ + Cu is 1.1V.\par
[Given : log 2 = 0.3, $\dfrac{2.303\ \mathrm{RT}}{\mathrm F}$ = 0.06]
\qn{5} How many of the following will not give iodoform test :\par
\qfig{p1_chem_s4_q05_1}
(g) CI$_3$$-$CHO \quad (h) PhCOPh
\qn{6} In how many reactions primary amine will be one of the products?\par
\qfig{p1_chem_s4_q06_1}
\raggedcolumns
\papersec{1}
\partsec{3}{MATHEMATICS}
\sect{SECTION-I (i)}
\qn{1} If $\displaystyle\sum_{i=1}^{10}(\mathrm x_i-5)=5$ and $\displaystyle\sum_{i=1}^{10}(\mathrm x_i-5)^2=125$, variance of x$_1$, x$_2$, \ldots\ldots x$_{10}$ will be -
\begin{qoptions}
\opt{A}{$\dfrac94$}
\opt{B}{$\dfrac{25}4$}
\opt{C}{$\dfrac{49}4$}
\opt{D}{$\dfrac{81}4$}
\end{qoptions}
\qn{2} Let A and B be two non-empty subsets of set X such that A is not a subset of B, then -
\begin{qoptions}
\opt{A}{A $\subseteq$ B$'$}
\opt{B}{B $\subseteq$ A$'$}
\opt{C}{A and B are disjoint}
\opt{D}{A and B$'$ are non disjoint}
\end{qoptions}
\qn{3} Let `R' be a relation over the set of natural numbers defined by (p, q) R(r, s) $\Rightarrow$ ps $-$ qr = 0. Then `R' is-
\begin{qoptions}
\opt{A}{Reflexive only}
\opt{B}{Symmetric only}
\opt{C}{Transitive only}
\opt{D}{Equivalence}
\end{qoptions}
\qn{4} Let B be a square matrix such that $|$B$|$ = 1 and 2A + B = B$^\mathrm T$ $-$ A$^\mathrm T$, then $|$A + B$|$ is -
\begin{qoptions}
\opt{A}{0}
\opt{B}{$-$1}
\opt{C}{1}
\opt{D}{2}
\end{qoptions}
\sect{SECTION-I (ii)}
\qn{5} If n is a positive integer and $\bigl(5+2\sqrt6\bigr)^{2n+1}$ = I + $f$, where I is an integer and 0 $\leq$ $f$ $<$ 1, then which of the following is/are true.
\begin{qoptions}
\opt{A}{I is an odd integer}
\opt{B}{I + 1 is divisible by 9}
\opt{C}{the integer next above $\bigl(5+2\sqrt6\bigr)^{2n+1}$ is divisible by 10}
\opt{D}{I $-$ 1 = $\dfrac{f}{1-f}$}
\end{qoptions}
\qn{6} If A and B are orthogonal matrices, then which of the following is/are also orthogonal matrix -
\begin{qoptions}
\opt{A}{A$^2$}
\opt{B}{AB}
\opt{C}{A$^2$B$^2$}
\opt{D}{A$^\mathrm T$B}
\end{qoptions}
\qn{7} If extremities of latus rectum of a hyperbola are (3, 5), (3, $-$1), ($-$1, 5), ($-$1, $-$1) and its transverse axis is parallel to x-axis, then -
\begin{qoptions}
\opt{A}{its eccentricity is 2}
\opt{B}{its eccentricity is greater than 2}
\opt{C}{one of its directrix is x = $\dfrac12$}
\opt{D}{one of its directrix is x = $\dfrac32$}
\end{qoptions}
\sect{SECTION-I (iii)}
\qn{8} India and Australia play a series of `n' one day matches and probability that India wins a match against Australia is $\dfrac12$.\par
{\footnotesize\setlength{\tabcolsep}{1.5pt}\noindent\begin{tabular}{|c|L{3.7cm}|c|c|}\hline & \multicolumn{1}{>{\centering\arraybackslash}p{3.7cm}|}{List-I} & & \multicolumn{1}{c|}{List-II} \\ \hline
(P) & If `n' is not fixed and series ends when any one of the team completes its 4$^\mathrm{th}$ win then probability that India wins the series is & (1) & $\dfrac4{2^7}$ \\ \hline
(Q) & If n = 7 then probability that India wins atleast three consecutive matches is & (2) & $\dfrac{47}{2^7}$ \\ \hline
(R) & For n = 7; probability that India wins series through consecutive wins is & (3) & $\dfrac12$ \\ \hline
(S) & For n = 7; probability that India doesn't win two consecutive matches is & (4) & $\dfrac{17}{2^6}$ \\ \hline
 & & (5) & $\dfrac{15}{2^6}$ \\ \hline
\end{tabular}}\par
\begin{qoptions}
\opt{A}{P $\to$ 4;Q $\to$ 2;R $\to$ 1;S $\to$ 3}
\opt{B}{P $\to$ 1;Q $\to$ 2;R $\to$ 3;S $\to$ 4}
\opt{C}{P $\to$ 4;Q $\to$ 3;R $\to$ 1;S $\to$ 2}
\opt{D}{P $\to$ 3;Q $\to$ 2;R $\to$ 1;S $\to$ 4}
\end{qoptions}
\qn{9} Match the statement in List-I with given value in List-II\par
{\footnotesize\setlength{\tabcolsep}{1.5pt}\noindent\begin{tabular}{|c|L{5.0cm}|c|c|}\hline & \multicolumn{1}{>{\centering\arraybackslash}p{5.0cm}|}{List-I} & & \multicolumn{1}{c|}{List-II} \\ \hline
(P) & Number of five digit numbers of the form d$_1$d$_2$d$_3$d$_4$d$_5$ where d$_\mathrm i$'s i = 1, 2, 3, 4, 5 are digits and satisfying d$_1$ $<$ d$_2$ $\leqslant$ d$_3$ $<$ d$_4$ $\leqslant$ d$_5$, is & (1) & $^{10}$C$_5$ \\ \hline
(Q) & Number of five digit numbers of the form d$_1$d$_2$d$_3$d$_4$d$_5$ where d$_\mathrm i$'s i = 1, 2, 3, 4, 5 are digits satisfying d$_1$ $>$ d$_2$ $\geqslant$ d$_3$ $>$ d$_4$ $>$ d$_5$ is & (2) & $^{11}$C$_4$ \\ \hline
(R) & Pankaj and Prateek play a unique game series in a chess tournament. They decide to play on till one of them wins 5 matches. If each match end only in win or loss. Number of ways in which series can be won by either of them, is & (3) & $^{11}$C$_6$ \\ \hline
(S) & A badminton team has to be selected comprising of 5 students out of 10 students for inter school tournament. Number of ways this can be done if a particular players is to be always included or always excluded from the team, is & (4) & 2.$^9$C$_5$ \\ \hline
 & & (5) & 5.$^9$C$_5$ \\ \hline
\end{tabular}}\par
\begin{qoptions}
\opt{A}{P $\to$ 3;Q $\to$ 3;R $\to$ 1,3;S $\to$ 1}
\opt{B}{P $\to$ 1,3;Q $\to$ 3,4;R $\to$ 2,3;S $\to$ 1}
\opt{C}{P $\to$ 3;Q $\to$ 1,3;R $\to$ 2;S $\to$ 4}
\opt{D}{P $\to$ 3;Q $\to$ 3;R $\to$ 1,4;S $\to$ 1,4}
\end{qoptions}
\qn{10} Match List-I with List-II and select the correct answer using the code given below the list.\par
{\footnotesize\setlength{\tabcolsep}{1.5pt}\noindent\begin{tabular}{|c|L{5.0cm}|c|c|}\hline & \multicolumn{1}{>{\centering\arraybackslash}p{5.0cm}|}{List-I} & & \multicolumn{1}{c|}{List-II} \\ \hline
(P) & If $\alpha$ + i$\beta$, $\alpha$,$\beta$ $\in$ R, $\beta$ $\neq$ 0 is a root of z$^5$ = 1, then value of $\dfrac{\left(\beta^4-\alpha^4\right)}{1-2\alpha^2}$ is & (1) & 0 \\ \hline
(Q) & If z$\bar{\mathrm z}$ = 1, then value of $\left|2+\dfrac1{\mathrm z}\right|^2$ + $|2-\mathrm z|^2$ & (2) & 1 \\ \hline
(R) & If $\omega$ = cos$\left(\dfrac{2\pi}7\right)$ + i sin $\dfrac{2\pi}7$ and $\alpha$ = $\omega$ + $\omega^2$ + $\omega^4$, $\beta$ = $\omega^3$ + $\omega^5$ + $\omega^6$, then value of 2$\alpha\beta$ is & (3) & 4 \\ \hline
(S) & If z$_1$, z$_2$, z$_3$, z$_4$ represent vertices of a parallelogram (in order), then z$_1$ $-$ z$_2$ + z$_3$ $-$ z$_4$ is & (4) & 10 \\ \hline
 & & (5) & 6 \\ \hline
\end{tabular}}\par
\begin{qoptions}
\opt{A}{P $\to$ 2;Q $\to$ 4;R $\to$ 3;S $\to$ 1}
\opt{B}{P $\to$ 2;Q $\to$ 4;R $\to$ 1;S $\to$ 3}
\opt{C}{P $\to$ 4;Q $\to$ 2;R $\to$ 3;S $\to$ 1}
\opt{D}{P $\to$ 3;Q $\to$ 1;R $\to$ 2;S $\to$ 4}
\end{qoptions}
\qn{11} Match List-I with List-II and select the correct answer using the code given below the list.\par
{\footnotesize\setlength{\tabcolsep}{1.5pt}\noindent\begin{tabular}{|c|L{5.0cm}|c|c|}\hline & \multicolumn{1}{>{\centering\arraybackslash}p{5.0cm}|}{List-I} & & \multicolumn{1}{c|}{List-II} \\ \hline
(P) & Let $\left(\vec a\times\vec b\right)\times\vec c=-5\vec a+4\vec b$ and $\vec a.\vec b=3$ and $\vec a\times\left(\vec b\times\vec c\right)=\lambda\vec b+\mu\vec c$, then $|\lambda|+|\mu|$ is\par (Where $\vec a$, $\vec b$, $\vec c$ are pairwise non-collinear) & (1) & 2 \\ \hline
(Q) & If the line y $-$ $\sqrt3$x + 3 = 0 cuts the parabola y$^2$ = x + 2 at A and B, given that $|$(PA)(PB)$|$ = $\dfrac{\lambda_1\left(2+\sqrt3\right)}{\lambda_2}$, then $\dfrac{\lambda_1^2+\lambda_2^2}5$ is\par (where P is $\left(\sqrt3,0\right)$ and $\lambda_1$, $\lambda_2$ are relatively prime) & (2) & 7 \\ \hline
(R) & The line $\dfrac{\mathrm x-4}1=\dfrac{\mathrm y-2}1=\dfrac{\mathrm z-\mathrm k}2$ lies on the plane 2x $-$ 4y + z = 9, then k is & (3) & 9 \\ \hline
(S) & Value of $\displaystyle\int_{-\pi/4}^{\pi/4}\dfrac{\mathrm x^9-3\mathrm x^7+\mathrm x^5+1}{\cos^2\mathrm x}$dx is equal to & (4) & 5 \\ \hline
 & & (5) & 4 \\ \hline
\end{tabular}}\par
\begin{qoptions}
\opt{A}{P $\to$ 1;Q $\to$ 2;R $\to$ 3;S $\to$ 4}
\opt{B}{P $\to$ 2;Q $\to$ 4;R $\to$ 3;S $\to$ 1}
\opt{C}{P $\to$ 4;Q $\to$ 1;R $\to$ 3;S $\to$ 2}
\opt{D}{P $\to$ 3;Q $\to$ 4;R $\to$ 2;S $\to$ 1}
\end{qoptions}
\sect{SECTION-II}
\qn{1} Number of ways in which 3 boys and 3 girls can be seated on a line where two particular girls do not want to sit adjacent to a particular boy is equal to -
\qn{2} A single die is rolled; then `n' coins are tossed, where n is the number shown on the die. If the probability of exactly four tails is k, then value of 384k is
\qn{3} The slope of tangent to the curve y = xe$^{\mathrm x^2}$ at the point of inflection is
\qn{4} Let $f$(x) = px$^3$ + qx$^2$ + rx have local maximum and minimum at x = 5 and x = 1 respectively. If $\displaystyle\int_{-1}^{1}f$(x)dx = 6, then value of $|$p + q + r$|$ is
\qn{5} If the solution of the differential equation y$^2$ + $\dfrac{\mathrm y}{1+\mathrm x^2}$ + y$'$tan$^{-1}$x + 2xyy$'$ = y$'$ (where y$'$ = $\dfrac{\mathrm{dy}}{\mathrm{dx}}$) is given by y(xy $-$ k + tan$^{-1}$x) = c (where c is constant of integration) then, k is equal to
\qn{6} Let $\alpha$, $\beta$, $\gamma$ and $\delta$ be the roots (real or non real) of equations x$^4$ $-$ 3x + 1 = 0. The value of $\alpha^3$ + $\beta^3$ + $\gamma^3$ + $\delta^3$ is
\raggedcolumns
\papersec{2}
\partsec{1}{PHYSICS}
\sect{SECTION-I (i)}
\qn{1} A uniform rope of mass m and length L hangs from the ceiling. A small wave pulse is created near bottom, Then which of following is correct regarding wave pulse ?\par
\qfig{p2_phy_s1_q01_1}
(i) Rate of change of speed with respect to time varies.\par
(ii) Rate of change of speed with respect to time is constant.\par
(iii) Rate of change of speed with respect to time is zero.\par
(iv) Wave speed varies
\begin{qoptions}
\opt{A}{i,iii}
\opt{B}{i,iv}
\opt{C}{ii, iv}
\opt{D}{iii, iv}
\end{qoptions}
\qn{2} The maximum velocity of an electron emitted by light of wavelength $\lambda$ incident on the surface of a metal of work function W is :-
\begin{qoptions}
\opt{A}{$\left[\dfrac{2(\mathrm{hc}+\lambda\mathrm W)}{\mathrm m\lambda}\right]^{1/2}$}
\opt{B}{$\dfrac{2(\mathrm{hc}-\lambda\mathrm W)}{\mathrm m}$}
\opt{C}{$\left[\dfrac{2(\mathrm{hc}-\lambda\mathrm W)}{\mathrm m\lambda}\right]^{1/2}$}
\opt{D}{$\left[\dfrac{2(\mathrm h\lambda-\mathrm W)}{\mathrm m}\right]^{1/2}$}
\end{qoptions}
\qn{3} An electron beam is accelerated by a potential difference V to hit a metallic target to produce X-rays. It produces continuous as well as characteristic X-rays. If $\lambda_\mathrm{min}$ is the smallest possible wavelength of X-ray in the spectrum, the variation of log $\lambda_\mathrm{min}$ with log V is correctly represented in :-
\begin{qoptions}
\opt{A}{\qfigi{p2_phy_s1_q03_A}}
\opt{B}{\qfigi{p2_phy_s1_q03_B}}
\opt{C}{\qfigi{p2_phy_s1_q03_C}}
\opt{D}{\qfigi{p2_phy_s1_q03_D}}
\end{qoptions}
\qn{4} If $\theta$ =110$^\circ$ then total number of images formed by the mirror system will be :-\par
\qfig{p2_phy_s1_q04_1}
\begin{qoptions}
\opt{A}{2}
\opt{B}{3}
\opt{C}{4}
\opt{D}{5}
\end{qoptions}
\sect{SECTION-I (ii)}
\qn{5} The p-T diagram of a cyclic process on an ideal gas is as shown.\par
\qfig{p2_phy_s2_q05_1}
Which of the following indicator diagrams is/are correctly represent the above cyclic process?
\begin{qoptions}
\opt{A}{\qfigi{p2_phy_s2_q05_A}}
\opt{B}{\qfigi{p2_phy_s2_q05_B}}
\opt{C}{\qfigi{p2_phy_s2_q05_C}}
\opt{D}{\qfigi{p2_phy_s2_q05_D}}
\end{qoptions}
\qn{6} A thin-walled, spherical conducting shell S of radius R is given charge Q. The same amount of charge is also placed at its centre C. Which of the following statement(s) are correct ?
\begin{qoptions}
\opt{A}{On the outer surface of S, the charge density is $\dfrac{\mathrm Q}{2\pi\mathrm R^2}$ .}
\opt{B}{The electric field is zero at all points inside S.}
\opt{C}{At a point just outside S, the electric field is double the field at a point just inside S.}
\opt{D}{At any point inside S, the electric field is inversely proportional to the square of its distance from C.}
\end{qoptions}
\qn{7} Switch S of the circuit shown in figure is closed at t = 0. If e denotes the induced e.m.f. in L and I is the current flowing through the circuit at time t, which of the following graphs is/are correct ?\par
\qfig{p2_phy_s2_q07_1}
\begin{qoptions}
\opt{A}{\qfigi{p2_phy_s2_q07_A}}
\opt{B}{\qfigi{p2_phy_s2_q07_B}}
\opt{C}{\qfigi{p2_phy_s2_q07_C}}
\opt{D}{\qfigi{p2_phy_s2_q07_D}}
\end{qoptions}
\sect{SECTION-II (i)}
\paraq{1}{2}
A semi cylinder made of a transparent plastic has a refraction index of n = $\sqrt2$ and a radius of R. There is a narrow incident light ray perpendicular to the flat side of the semi cylinder at d distance from the axis of symmetry.\par
\qfig{p2_phy_s3_para1_2_1}
\qn{1} What can the maximum value of d (in cm) be so that the light ray can still leave the other side of the semi cylinder? [Take : R = 10 cm]
\qn{2} When the value of d chosen is such that TIR just takes place then time (in nanoseconds) for which light remains inside the cylinder is (take R = 9cm \& C = 3 $\times$ 10$^8$ m/s)
\paraq{3}{4}
In the given circuit all capacitor were uncharged initially. First switch S$_1$ is closed for long time when circuit reaches to steady state, switch S$_1$ is opened \& S$_2$ is closed.\par
\qfig{p2_phy_s3_para3_4_1}
\qn{3} The current through R$_2$, just after closing the switch S$_2$ is (in A) :-
\qn{4} Net heat produced in R$_2$ after closing switch S$_2$ (in $\mu$J) :-
\sect{SECTION-II (ii)}
\qn{5} In the figure shown the current I in ampere is :- (R = 1$\Omega$)\par
\qfig{p2_phy_s4_q05_1}
\qn{6} A square of side L/2 is removed from one corner of a square sandwich that has sides of length L. The center of mass of the remainder of the sandwich moves from C to C$'$. The displacement of the center of mass (from C to C$'$) is p$\dfrac{\sqrt2}{12}$L. The value of p is :-\par
\qfig{p2_phy_s4_q06_1}
\qn{7} In an $\alpha$-decay the kinetic energy of $\alpha$ particle is 48 MeV and Q value of the reaction is 50 MeV. The mass number of the parent nucleus is 20n then find n : (Assume that daughter nucleus is in ground state)
\qn{8} The potential energy of a particle of mass m is given by V (x) = $\begin{cases}\mathrm E_0&0\leqslant \mathrm x\leqslant 1\\0&\mathrm x>1\end{cases}$ $\lambda_1$ and $\lambda_2$ are the de-Broglie wavelengths of the particle, when 0 $\leq$ x $\leq$ 1 and x $>$ 1 respectively. If the total energy of particle is 2E$_0$, find $(\lambda_1/\lambda_2)^2$.
\qn{9} A tuning fork of frequency 480 Hz is used in an experiment for measuring speed of sound ($\nu$) in air by resonance tube method. Resonance is observed to occur at two successive lengths of the air column, $\ell_1$ = 30 cm and $\ell_2$ = 70 cm. Then $\nu$ (in m/s) is equal to:
\qn{10} In a Young's double slit experiment 15 fringes are observed on a small portion of the screen when light of wavelength 500 nm is used. Ten fringes are observed on the same section of the screen when another light source of wavelength $\lambda$ is used. Then the value of $\lambda$ (in nm) \underline{\hspace{1cm}}.
\raggedcolumns
\papersec{2}
\partsec{2}{CHEMISTRY}
\sect{SECTION-I (i)}
\qn{1} Colligative properties of the solution depend upon :
\begin{qoptions}
\opt{A}{Nature of the solution}
\opt{B}{Nature of the solvent}
\opt{C}{Number of solute particles}
\opt{D}{Number of moles of solvent}
\end{qoptions}
\qn{2} Correct formula of apatite family is :\par
(where X = F, Cl , OH)
\begin{qoptions}
\opt{A}{Ca$_3$(PO$_4$)$_2$.CaX$_2$}
\opt{B}{2Ca$_3$(PO$_4$)$_2$.CaX$_2$}
\opt{C}{Ca$_3$(PO$_4$)$_2$.2CaX$_2$}
\opt{D}{3Ca$_3$(PO$_4$)$_2$.CaX$_2$}
\end{qoptions}
\qn{3} \qfigi{p2_chem_s1_q03_R}
\begin{qoptions}
\opt{A}{\qfigi{p2_chem_s1_q03_A2}}
\opt{B}{\qfigi{p2_chem_s1_q03_B2}}
\opt{C}{\qfigi{p2_chem_s1_q03_C2}}
\opt{D}{\qfigi{p2_chem_s1_q03_D2}}
\end{qoptions}
\qn{4} \qfigi{p2_chem_s1_q04_R}\par
Correct statement(s) for above reaction is :
\begin{qoptions}
\opt{A}{(C) is cyclohexanol}
\opt{B}{(C) is \qfigi{p2_chem_s1_q04_B2}}
\opt{C}{(C) is benzene}
\opt{D}{(C) is \qfigi{p2_chem_s1_q04_D2}}
\end{qoptions}
\sect{SECTION-I (ii)}
\qn{5} Three different solutions of oxidising agents KMnO$_4$, K$_2$Cr$_2$O$_7$ and I$_2$ is titrated separately with 0.158 gm of Na$_2$S$_2$O$_3$. If molarity of each oxidising agent is 0.1 M and reactions are :\par
MnO$_4^-$ + S$_2$O$_3^{2-}$ $\longrightarrow$ MnO$_2$ + SO$_4^{2-}$\par
Cr$_2$O$_7^{2-}$ + S$_2$O$_3^{2-}$ $\longrightarrow$ Cr$^{3+}$ + SO$_4^{2-}$\par
I$_2$ + S$_2$O$_3^{2-}$ $\longrightarrow$ S$_4$O$_6^{2-}$ + I$^-$\par
Then,
\begin{qoptions}
\opt{A}{Volume of KMnO$_4$ used is maximum}
\opt{B}{Volume of iodine used is minimum}
\opt{C}{Wt. of I$_2$ used in titration is maximum}
\opt{D}{Gram equivalent of Na$_2$S$_2$O$_3$ are same in all the reactions.}
\end{qoptions}
\qn{6} The cubic unit cell structure of a compound containing cation M and anion X is shown below. When compared to the anion, the cation has smaller ionic radius. Choose the correct statement(s) :\par
\qfig{p2_chem_s2_q06_X}
\begin{qoptions}
\opt{A}{The empirical formula of the compound is MX.}
\opt{B}{The cation M and anion X have different coordination geometries.}
\opt{C}{The ratio of M-X bond length to the cubic unit cell edge length is 0.866.}
\opt{D}{The ratio of the ionic radii of cation M to anion X is 0.414.}
\end{qoptions}
\qn{7} Which of following are ores of Pb?
\begin{qoptions}
\opt{A}{Calamine}
\opt{B}{Galena}
\opt{C}{Cerrusite}
\opt{D}{Cassiterite}
\end{qoptions}
\sect{SECTION-II (i)}
\paraq{1}{2}
0.20 moles of CO taken in a 2.463 litre flask is maintained at 750K along with a catalyst so that the following reaction can take place: CO(g) + 2H$_2$(g) $\rightleftharpoons$ CH$_3$OH(g). Hydrogen is introduced until the total pressure of the system is 7.5 atm at equilibrium and 0.1 mole of methanol is formed.\par
(R = 0.0821 L-atm/K-mol).
\qn{1} For the reaction, K$_\mathrm P$ (in atm$^{-2}$) is
\qn{2} What is the final pressure (in atm) if the same amount of CO and H$_2$ as before used, but with no catalyst so that the reaction does not take place?
\paraq{3}{4}
KMnO$_4$ act as oxidising agent in acidic, basic and neutral medium.\par
In acidic medium KMnO$_4$ oxidises I$^-$ to I$_2$ , C$_2$O$_4^{2-}$ into CO$_2$\par
Answer the following questions.
\qn{3} Number of moles of KMnO$_4$ needed to oxidises 3 moles of C$_2$O$_4^{2-}$ in acidic medium is ........
\qn{4} In alkaline medium KMnO$_4$ oxidises KI into a potassium salt `X' and itself is converted into Y. Sum of atomicity of X and oxidation number of metal in Y is ....
\sect{SECTION-II (ii)}
\qn{5} At 298 K, pure T$_2$O has pT (like pH) is 7.60. Find out the pT of a solution prepared by adding 10 ml of 0.2 M TCl to 15 ml of 0.20 M NaOT. (Given log 2 = 0.30)\par
Report your answer by multiplying actual answer with 100.
\qn{6} Consider following statements :\par
(I) CrO is basic oxide\par
(II) Mn$_2$O$_7$ is green oil\par
(III) K$_2$Cr$_2$O$_7$ act as oxidising agent in basic medium\par
(IV) Stability Bi$^{3+}$ $>$ Bi$^{5+}$\par
(V) Aqua regia is mixture of conc. HCl and conc. HNO$_3$ in 3 : 1 molar\par
(VI) Common oxidation state of lanthanides is +3\par
(VII) Highest oxidation state shown by actinides is +7\par
Number of correct statements is.
\qn{7} Number of correct statements among following is.\par
(I) Barium salt imparts apple green colour to oxidising flame\par
(II) In Mac-Arthur Forest cyanide process Zn act as reducing agent\par
(III) K$_2$CO$_3$ cannot be prepared by Solvay soda process\par
(IV) Fe(CO)$_5$ is diamagnetic and having trigonal bipyramidal geometry\par
(V) [Co(H$_2$O)$_3$F$_3$] is high spin complex\par
(VI) EA$_1$ : F $>$ Cl $>$ Br $>$ I\par
(VII) Bond energy : Cl$_2$ $>$ F$_2$ $>$ Br$_2$ $>$ I$_2$
\qn{8} Which of the following reaction(s) can produce \qfigi{p2_chem_s4_q08_T} as one of the product ?\par
(i) \qfigi{p2_chem_s4_q08_i}\par
(ii) \qfigi{p2_chem_s4_q08_ii}\par
(iii) \qfigi{p2_chem_s4_q08_iii}\par
(iv) \qfigi{p2_chem_s4_q08_iv}\par
(v) \qfigi{p2_chem_s4_q08_v}
\qn{9} Which of the following reaction is showing correct major product :\par
(i) \qfigi{p2_chem_s4_q09_i}\par
(ii) \qfigi{p2_chem_s4_q09_ii}\par
(iii) \qfigi{p2_chem_s4_q09_iii}\par
(iv) \qfigi{p2_chem_s4_q09_iv}\par
(v) \qfigi{p2_chem_s4_q09_v}
\qn{10} In how many of the following reactions phenol will be one of the products?\par
(i) \qfigi{p2_chem_s4_q10_i}\par
(ii) \qfigi{p2_chem_s4_q10_ii}\par
(iii) \qfigi{p2_chem_s4_q10_iii}\par
(iv) \qfigi{p2_chem_s4_q10_iv}\par
(v) \qfigi{p2_chem_s4_q10_v}\par
(vi) \qfigi{p2_chem_s4_q10_vi}
\raggedcolumns
\papersec{2}
\partsec{3}{MATHEMATICS}
\sect{SECTION-I (i)}
\qn{1} If z$_1$, z$_2$ \& z$_3$ are three complex numbers with modulus 2, 3 \& 4 respectively and if $|\mathrm z_1+\mathrm z_2+\mathrm z_3|$ = 5, then $|4\mathrm z_2\mathrm z_3+9\mathrm z_3\mathrm z_1+16\mathrm z_1\mathrm z_2|$ is -
\begin{qoptions}
\opt{A}{24}
\opt{B}{60}
\opt{C}{120}
\opt{D}{240}
\end{qoptions}
\qn{2} If z is a complex number such that z + i$|$z$|$ = i$\bar{\mathrm z}$ + 1, then $|$z$|$ is -
\begin{qoptions}
\opt{A}{4}
\opt{B}{2}
\opt{C}{1}
\opt{D}{0}
\end{qoptions}
\qn{3} The x-intercept of the plane which is passing through the intersection of the planes\par
x + 2y + 3z + 5 = 0 and 2x $-$ 3y + 7z + 1 = 0 and parallel to line $\vec{\mathrm r}$ = $\hat{\mathrm i}$ + 2$\hat{\mathrm j}$ + $\lambda\left(8\hat{\mathrm i}-7\hat{\mathrm j}-4\hat{\mathrm k}\right)$, is -
\begin{qoptions}
\opt{A}{$\dfrac{13}5$}
\opt{B}{$-\dfrac75$}
\opt{C}{$\dfrac65$}
\opt{D}{$-\dfrac95$}
\end{qoptions}
\qn{4} Equation of the plane passing through line of intersection of the planes 2x $-$ 4y + 3z = 5 and x + 2y + z = 4 and perpendicular to the plane 2x $-$ y + z = 4 is -
\begin{qoptions}
\opt{A}{$-$9x + 26y + 8z + 39 = 0}
\opt{B}{9x + 26y + 8z $-$ 39 = 0}
\opt{C}{9x $-$ 26y $-$ 8z $-$ 39 = 0}
\opt{D}{9x + 26y + 8z + 39 = 0}
\end{qoptions}
\sect{SECTION-I (ii)}
\qn{5} If A = $\begin{bmatrix}1&2\\2&1\end{bmatrix}$ \& A$^\mathrm n$ = $\begin{bmatrix}\mathrm a&\mathrm b\\\mathrm c&\mathrm d\end{bmatrix}$, (where n $\geq$ 2 \& n $\in$ N), then -
\begin{qoptions}
\opt{A}{a = d}
\opt{B}{b = c}
\opt{C}{b = a + 1 if n is odd}
\opt{D}{b = a $-$ 1 if n is even}
\end{qoptions}
\qn{6} $\displaystyle\lim_{\mathrm n\to\infty}\sum_{\mathrm i=1}^{\mathrm n}\dfrac{\mathrm i}{\mathrm n^2+\mathrm i^2}$ is equal to -
\begin{qoptions}
\opt{A}{$\dfrac14-\ell\mathrm n2$}
\opt{B}{$\dfrac12-\ell\mathrm n2$}
\opt{C}{$\ell\mathrm n\sqrt2$}
\opt{D}{$\dfrac12\ell\mathrm n2$}
\end{qoptions}
\qn{7} System of equation x + y + az = b, 2x + 3y = 2a \& 3x + 4y + a$^2$z = ab + 2 has
\begin{qoptions}
\opt{A}{unique solution when a $\neq$ 0, b $\in$ R}
\opt{B}{no solution when a = 0, b = 1}
\opt{C}{infinite solution when a = 0, b = 2}
\opt{D}{infinite solution when a = 1, b $\in$ R}
\end{qoptions}
\sect{SECTION-II (i)}
\paraq{1}{2}
Let bag `A' contains 2 Red \& 3 Green balls \& bag `B' contains 3 Red \& 2 Green balls. If two balls are drawn randomly from each of bag A \& B and then put in empty bag C.\par
\qn{1} Probability that bag C contains 3 Red balls, is -
\qn{2} If it is found that bag C contains equal number of Red and Green balls, then probability that 2 Red balls were drawn from bag B, is -
\paraq{3}{4}
There are 5 boys and 5 girls.\par
\qn{3} If they are seated in a straight line, then the number of ways in which all boys and all girls are together is $\lambda$(5!5!), then the value of $\lambda$ is -
\qn{4} Number of ways in which they can be divided into two groups, if first group contains 2 boys \& 3 girls, second group contains 3 boys and 2 girls is -
\sect{SECTION-II (ii)}
\qn{5} Consider the curves C$_1$ : $|$z $-$ 2$|$ = 2 + Re(z) and C$_2$ : $|$z$|$ = 3 (where z = x + iy, x, y $\in$ R and i = $\sqrt{-1}$). They intersect at P and Q in the first and fourth quadrants respectively. Tangents to C$_1$ at P and Q intersect the x-axis at R and tangents to C$_2$ at P and Q intersect the x-axis at S. If area of $\Delta$PRS is $\lambda\sqrt2$ sq.units, then ($\lambda^2$) is
\qn{6} A plane passes point P(0, 1, $-$1) and contains a line L : $\dfrac{\mathrm x-1}2=\dfrac{\mathrm y-1}3=\dfrac{\mathrm z+2}1$. If R is a point on line L such that $\Delta$PQR is right angled at P, where Q(1, 1, $-$2), then the value of [PR] is (where [.] denotes greatest integer function)
\qn{7} If the ratio of 2$^\mathrm{nd}$ to 3$^\mathrm{rd}$ term in the expansion of (x + y)$^\mathrm n$ is equal to the ratio of 3$^\mathrm{rd}$ to 4$^\mathrm{th}$ term in the expansion of (x + y)$^{\mathrm n+3}$, then n is equal to
\qn{8} If the function $f$(x) = sin$^3$x + msin$^2$x has exactly one point of local minima and one point of local maxima in x $\in$ $\left(-\dfrac\pi2,\dfrac\pi2\right)$, then number of possible integral values of `m' is
\qn{9} Images of circle C : x$^2$ + y$^2$ $-$ 8x $-$ 10y + 16 = 0 are taken in variable line ($\lambda$ + 1)x + (1 $-$2$\lambda$)y + $\lambda$ = 2 $\forall$ $\lambda$ $\in$ R. If locus of centre of there images intersects circle `C' at A and B, then $\dfrac{\mathrm{AB}}{\sqrt3}$, is
\qn{10} If the vectors $\overrightarrow{\mathrm{PQ}}$ = $-$4$\hat{\mathrm i}$ + 5$\hat{\mathrm k}$ and $\overrightarrow{\mathrm{PR}}$ = 3$\hat{\mathrm i}$ $-$ $\lambda\hat{\mathrm j}$ + $\hat{\mathrm k}$ (where $\lambda$ $>$ 0 are the sides of a $\Delta$PQR). Length of median through P is $\sqrt{24}$ and length of the side QR is $\sqrt{\mathrm k}$, then sum of digits of k is
\end{multicols}
\end{document}
